% 题证摘录；来源与整理边界见相同编号分题文件。

\paragraph{Situation 33.47.2.}
Let $k$ be a field, let $X$ be a scheme over $k$, let $\mathcal L$ be an invertible $\mathcal O_X$-module, let $V$ be a finite dimensional $k$-vector space, and let $\psi:V\to\Gamma(X,\mathcal L)$ be a $k$-linear map. Say $\dim(V)=r$ and we have a basis $v_1,\ldots,v_r$ of $V$. Then we obtain a ``universal divisor''
\[H_{univ}=Z(s_{univ})\subset\mathbf A^r\times_k X\]
as the zero scheme (Divisors, Definition 31.15.8) of the section
\[s_{univ}=\sum_{i=1,\ldots,r}x_i\psi(v_i)\in\Gamma(\mathbf A^r\times_k X,\operatorname{pr}_2^*\mathcal L).\]
For a field extension $k'/k$ the $k'$-points $v\in\mathbf A^r_k(k')$ correspond to vectors $(a_1,\ldots,a_r)$ of elements of $k'$. Thus we may on the one hand think of $v$ as the element $v=\sum_{i=1,\ldots,r}a_iv_i\in V\otimes_k k'$ and on the other hand we may assign to $v$ the section
\[\psi(v)=\sum_{i=1,\ldots,r}a_i\psi(v_i)\in\Gamma(X_{k'},\mathcal L|_{X_{k'}}).\]
With this notation it is clear that the fibre of $H_{univ}$ over $v\in V\otimes k'$ is the zero scheme of $\psi(v)$. In a formula:
\[H_v=H_{univ,v}=Z(\psi(v)).\]
We will denote this common value by $H_v$ as indicated. Finally, in this situation let $P$ be a property of vectors $v\in V\otimes_k k'$ for $k'/k$ an arbitrary field extension. We say $P$ holds for general $v\in V\otimes_k k'$ if there exists a nonempty Zariski open $U\subset\mathbf A^r_k$ such that if $v$ corresponds to a $k'$-point of $U$ for any $k'/k$ then $P(v)$ holds.

\begin{lemma}[33.47.3]
In Situation 33.47.2 assume
\begin{enumerate}
\item $X$ is smooth over $k$,
\item the image of $\psi:V\to\Gamma(X,\mathcal L)$ generates $\mathcal L$,
\item the corresponding morphism $\varphi_{\mathcal L,\psi}:X\to\mathbf P(V)$ is an immersion.
\end{enumerate}
Then for general $v\in V\otimes_k k'$ the scheme $H_v$ is smooth over $k'$.
\end{lemma}
\begin{proof}
(We observe that $X$ is separated and finite type as a locally closed subscheme of a projective space.) Let us use the notation introduced above the statement of the lemma. We consider the projections
\[
\xymatrix{\mathbf A^r_k\times_k X\ar[d]&H_{univ}\ar[l]\ar[ld]^p\ar[r]\ar[rd]_q&\mathbf A^r_k\times_k X\ar[d]\\X&&\mathbf A^r_k}
\]
Let $\Sigma\subset H_{univ}$ be the singular locus of the morphsm $q:H_{univ}\to\mathbf A^r_k$, i.e., the set of points where $q$ is not smooth. Then $\Sigma$ is closed because the smooth locus of a morphism is open by definition. Since the fibre of a smooth morphism is smooth, it suffices to prove $q(\Sigma)$ is contained in a proper closed subset of $\mathbf A^r_k$. Since $\Sigma$ (with reduced induced scheme structure) is a finite type scheme over $k$ it suffices to prove $\dim(\Sigma)<r$. This follows from Lemma 33.20.4. Since dimensions aren't changed by replacing $k$ by a bigger field (Morphisms, Lemma 29.29.3), we may and do assume $k$ is algebraically closed.

By dimension theory (Lemma 33.20.4), it suffices to prove that for $x\in X\setminus Z$ closed we have $p^{-1}(\{x\})\cap\Sigma$ has dimension $<r-\dim(X')$ where $X'$ is the unique irreducible component of $X$ containing $x$. As $X$ is smooth over $k$ and $x$ is a closed point we have $\dim(X')=\dim\mathfrak m_x/\mathfrak m_x^2$ (Morphisms, Lemma 29.35.12 and Algebra, Lemma 10.140.1). Thus we win if
\[\dim p^{-1}(x)\cap\Sigma<r-\dim\mathfrak m_x/\mathfrak m_x^2\]
for all $x\in X$ closed.

Since $V$ globally generated $\mathcal L$, for every irreducible component $X'$ of $X$ there is a nonempty Zariski open of $\mathbf A^r$ such that the fibres of $q$ over this open do not contain $X'$. (For example, if $x'\in X'$ is a closed point, then we can take the open corresponding to those vectors $v\in V$ such that $\psi(v)$ does not vanish at $x'$. This open will be the complement of a hyperplane in $\mathbf A^r_k$.) Let $U\subset\mathbf A^r$ be the (nonempty) intersection of these opens. Then the fibres of $q^{-1}(U)\to U$ are effective Cartier divisors on the fibres of $U\times_k X\to U$ (because a nonvanishing section of an invertible module on an integral scheme is a regular section). Hence the morphism $q^{-1}(U)\to U$ is flat by Divisors, Lemma 31.19.9. Thus for $x\in X$ closed and $v\in V=\mathbf A^r_k(k)$, if $(x,v)\in H_{univ}$, i.e., if $x\in H_v$ then $q$ is smooth at $(x,v)$ if and only if the fibre $H_v$ is smooth at $x$, see Morphisms, Lemma 29.35.14.

Consider the image $\psi(v)_x$ in the stalk $\mathcal L_x$ of the section corresponding to $v\in V$. We have
\[x\in H_v\ \Leftrightarrow\ \psi(v)_x\in\mathfrak m_x\mathcal L_x.\]
If this is true, then we have
\[H_v\text{ singular at }x\ \Leftrightarrow\ \psi(v)_x\in\mathfrak m_x^2\mathcal L_x.\]
Namely, $\psi(v)_x$ is not contained in $\mathfrak m_x^2\mathcal L_x$ $\Leftrightarrow$ the local equation for $H_v\subset X$ at $x$ is not contained in $\mathfrak m_x^2$ $\Leftrightarrow$ $\mathcal O_{H_v,x}$ is regular (Algebra, Lemma 10.106.3) $\Leftrightarrow$ $H_v$ is smooth at $x$ over $k$ (Algebra, Lemma 10.140.5). We conclude that the closed points of $p^{-1}(x)\cap\Sigma$ correspond to those $v\in V$ such that $\psi(v)_x\in\mathfrak m_x^2\mathcal L_x$. However, as $\varphi_{\mathcal L,\psi}$ is an immersion the map
\[V\longrightarrow\mathcal L_x/\mathfrak m_x^2\mathcal L_x\]
is surjective (small detail omitted). By the above, the closed points of the locus $p^{-1}(x)\cap\Sigma$ viewed as a subspace of $V$ is the kernel of this map and hence has dimension $r-\dim\mathfrak m_x/\mathfrak m_x^2-1$ as desired.
\end{proof}
